\subsection{One normalization for the fourth theta tensor}
\label{subsec:theta-normalization}

Local integral models can have different generators. The fourth tensor needs a scalar normalization compatible with all the translates occurring in one period.

Fix an imaginary quadratic field \(K=\mathbb Q(\sqrt{-d})\)
with \(d\) prime, \(d\equiv7\bmod8\). Put
\(\tau=(b+\sqrt{-d})/2\) and take
\(\Gamma=\lambda(a\mathbb Z+\tau\mathbb Z)\), where
\(4ac=b^2+d\), \(a>0\) is odd and \(\lambda\ne0\).
Thus \(b\) is odd and \(c\) is even. Write
\(H_0(z,w)=\sum z_i\overline{w_i}/\operatorname{area}(\Gamma)\).

For a positive Hermitian form \(H\), linear in the first
variable, put \(E=\Im H\). A CM lattice \(M\) here is a full
\(\mathcal O_K\)-lattice self-dual for \(E\). Its sign character is
\[
  \alpha(x)=\exp\big(\pi i\,\sqrt d\,H(x,x)/2\big).
\]
Indeed the quadratic expression in the exponent divided by
\(\pi i\) is \(E(\tau x,x)\). For \(x,y\in M\) its
polarization is \(2E(\tau x,y)-bE(x,y)\), congruent to
\(E(x,y)\bmod2\). Consequently

\[
  \alpha(x+y)=\alpha(x)\alpha(y)e^{\pi i E(x,y)}.
\]
On one copy of \(\Gamma\) the quadratic expression is
\(au^2+buv+cv^2\) at \(\lambda(au+\tau v)\). Its parity is
\(u+uv\), the sign of the odd sigma section translated by
the half period \(\lambda a/2\). In an orthogonal sum the
signs multiply. Let \(f_\nu\) denote the translated one-variable sigma section in the \(\nu\)-th coordinate, divided by its nonzero value at zero. The product section in these coordinates is
\[
  F_0((z_\nu)_\nu)=\prod_\nu f_\nu(z_\nu),\qquad f_\nu(0)=1,\qquad F_0(0)=1.
\]

Take a positive CM Hermitian space of rank \(n\) with such
a lattice \(M\), in the complex scale of \(\Gamma\).
There is a rational isometry \(g_0\) of its fourfold sum to
the standard space of rank \(m=4n\). In fact diagonalize the
Hermitian form over \(K\) relative to the rational CM form
underlying \(H_0\). Every diagonal entry is a positive rational
number \(r\).

For a positive entry \(r=a/b\) with \(a,b\in\mathbb Z_{>0}\), Lagrange's theorem gives
\[
  ab=n_1^2+\cdots+n_4^2,\qquad r=\sum_{i=1}^4(n_i/b)^2.
\]
 The matrix of multiplication by the resulting
quaternion has transpose times itself equal to \(r I_4\).
Taking the block diagonal sum of the four-dimensional
quaternion matrices gives \(g_0\). Its equality of forms is
\[
       g_0^*H_0^{\oplus 4n}g_0=H^{\oplus4}
\]
in the corresponding ordered coordinates, with \(H_0\)
the standard form on one line. In particular \(L_0=g_0 M^4\) and \(S=\Gamma^m\)
are commensurable in the common rational space.

\paragraph{Choosing a local isometry and its determinant.}

To transport the normalized theta tensor from the standard lattice, we need a local isometry with the prescribed determinant.

\begin{lemma}[Local comparison with the standard CM lattice]
\label{lem:cm-standard-local-isometry}
For each finite prime \(p\) the lattices \(L_0\) and \(S\) are
locally isometric. The isometry at \(p\) can be chosen rational
with \(1-\delta\) invertible. For \(L_1=\delta L_0\), its intersection \(I=L_1\cap S\) satisfies
\[
 p\nmid[L_1:I]\,[S:I].
\]

\end{lemma}

\begin{proof}
After removing the common complex scale, write
\(\mathfrak a=a\mathbb Z+\tau\mathbb Z\) for the ideal of the
reference curve. It is invertible with

\[
  \mathfrak a\overline{\mathfrak a}=a\mathcal O_K,\qquad {\rm N}\mathfrak a=a
\]
 since
\(\mathrm N\mathfrak a=a\).

Tracing the Hermitian pairing
\(h(x,y)=\sum x_i\bar y_i/a\) against
\(1/\sqrt{-d}\) gives \(E(x,y)\). Since the inverse
different is \((1/\sqrt{-d})\mathcal O_K\), stability under
\(\mathcal O_K\) identifies self-duality for \(E\) with
unimodularity for \(h\). Locally choose a generator of
\(\mathfrak a\), so \(S\) has Gram matrix a unit times the
identity. At a split prime the two direct summands of a lattice
are in perfect duality. For a basis \(e_i\) of one summand and the dual basis \(e'_j\) of the other, the split pairing has entries
\[
  h(e_i,e'_j)=\delta_{ij}.
\]
This description includes \(p=2\).

At an inert odd prime, reducing a
unimodular form gives a nondegenerate Hermitian form over
\(\mathbb F_{p^2}\). It has a vector of nonzero norm by
polarization. Lift the vector, split its orthogonal complement
and iterate. The norm on unramified units is surjective, by
finite-field norm surjectivity and lifting the norm using the
surjective residue trace.
For the successive corrections to a norm, the linear term is the residue trace: for \(n\ge1\),
\[
  N(1+p^nx)\equiv 1+p^n\operatorname{Tr}(x)\pmod {p^{n+1}}.
\]
 Thus all diagonal units may be matched.

At the odd ramified prime the residual forms are nondegenerate
symmetric forms. They have the same determinant square class
because the generic change of basis has determinant of norm a
unit, and the residue of that norm is a square. Dimension and
determinant classify nonsingular forms over the finite field of
odd order. Lift a residual isometry modulo the uniformizer.
For a resulting Gram matrix \(B\) close to the desired \(C\),
\(B^{-1}C\) is self-adjoint for \(B\) and congruent to one.
Put \(A=B^{-1}C\) and \(R=A^{1/2}\).
The binomial series is integral since \(A\) is congruent
to \(1\) and \(1/2\in\mathbb Z_p\). Each power is
self-adjoint for \(B\), hence so is \(R\), and
\[
                  R^* B R=BR^2=BA=C .
\] This proves the local assertion.

The transporter of the two lattices is open in the local
unitary group. It contains elements for which both \(1+\delta\)
and \(1-\delta\) are invertible. For example multiply a
transporter by nearby scalar norm-one elements, avoiding its
finitely many forbidden eigenvalues. On this open subset
the Cayley change \(X=(1-\delta)(1+\delta)^{-1}\) parametrizes
the group by skew-Hermitian matrices. The rational skew-Hermitian
matrices are dense in their localization since they form a
rational vector space. We can therefore take \(\delta\) rational,
with \(\delta L_{0,p}=S_p\) and \(1-\delta\) invertible.
Put \(g_1=\delta g_0\). Each intersection index of
\(L_1=g_1M^4\) and \(S\) is prime to \(p\).
\end{proof}

For the choices in Lemma~\ref{lem:cm-standard-local-isometry}
one algebraic number suffices for normalization. Let
\(J=(\bigwedge^m L_0)/(\bigwedge^m S)\) be the fractional
ideal ratio in the common determinant line, with complex scale
removed. At every prime the local transporter has norm-one
determinant and \((\det\delta)J_p=\mathcal O_{K,p}\).
Hence \(J\bar J=\mathcal O_K\). Choose \(h>0\) with
\(J^h=(u)\), using finiteness of the ideal class group.
The positive rational number \(u\bar u\) generates the unit
ideal, so equals one. Fix any \(\kappa^{2h}=u\). All complex
conjugates of \(\kappa\) have modulus one. For every embedding
at \(p\) and every such \(\delta\) we have
\begin{equation}\label{eq:kappa-determinant}
             |\kappa|_p^2 |\det\delta|_p=1 .
\end{equation}
In particular the two valuations above a split prime are
accounted for separately.

\begin{lemma}[Integral Cayley change of the theta isometry]\label{lem:theta-cayley-scalar}
For the self-dual CM lattices with the intrinsic sign above,
use an isometry \(\delta\) for which the stated Cayley
operators are defined. The transport and its scalar satisfy
\[
 F\circ\delta^{-1}=c_\delta T_{L,\delta L}F,\qquad c_\delta^8=(\det\delta)^4.
\]

\end{lemma}

\begin{proof}
The theta line of a self-dual lattice with the given
semicharacter consists of the entire functions fixed by
\(\alpha(x)U_x\) for every lattice vector, where
\[
 U_xF(z)=e^{-\pi H(z,x)-\pi H(x,x)/2}F(z+x).
\]
It has dimension one by Riemann--Roch for a principal
polarization: the positive line bundle has vanishing higher
cohomology and top self-intersection divided by \(m!\) equal
to the Pfaffian of \(E\), of absolute value one. Equivalently
one can choose a symplectic basis and solve the automorphy
relations for its Fourier coefficients. For commensurable
self-dual \(L_i,L_j\) with the same quadratic sign on their
intersection, define
\[
 T_{ij}F=[L_j:L_i\cap L_j]^{-1/2}
       \sum_{x\in L_j/(L_i\cap L_j)}\alpha(x)U_xF .
\]
The Heisenberg multiplication rule

\[
  U_x U_y=e^{-\pi i E(x,y)}U_{x+y}
\]
 shows independence of
representatives and the fixed-vector property on \(L_j\).
Use the Gaussian metric and normalized Haar average on any
common period domain. The distinct summands are orthogonal:
translation by \(\ell\in L_i\) distinguishes the characters
\(e^{2\pi i E(x,\ell)}\), by self-duality of \(L_i\).
Each has the norm of \(F\), since the translated density has
the same mean. Thus \(T_{ij}\) is an isometry of lines.

For three such lattices carrying the intrinsic quadratic signs,
\(T_{23}T_{12}=\xi T_{13}\) with \(\xi^8=1\). To compute the
phase, project onto the initial line by averaging its
translations, or equivalently retain the coefficient of the
trivial \(L_1\)-character in the preceding orthogonal expansion.
Up to positive index factors the scalar is the sum over pairs
\(z\in L_3,y\in L_2,z+y\in L_1\) of
\[
 \alpha(z)\alpha(y)\alpha(z+y)e^{-\pi i E(z,y)} .
\]
For a common finite quotient take \(z,y\) modulo
\(2(L_1\cap L_2\cap L_3)\) in their respective lattices.
This counts each contributing pair of cosets in the two
intertwiners with the same multiplicity: translates of
\(z\) by \(t\in L_2\cap L_3\) are paired with translates of
\(y\) by \(-t\), and \(y\) can independently be shifted by
\(L_1\cap L_2\).

The Heisenberg law leaves the summand
unchanged under both operations. The displayed function
on the finite group of pairs is a homogeneous quadratic
character, by its quadratic exponent. Its sum is nonzero
because the intertwining maps are isometries and the
coefficient of the initial line in \(T_{13}\) is positive.
On the radical of its polarization its restriction is a
character.
Indeed on the radical \(R\) the defining polarization is trivial, so
\[
  q(x+r)=q(x)q(r)\qquad(r\in R).
\]
 A nontrivial such restriction would annihilate
the sum by translation.
In fact translation by a radical element \(r\) multiplies the sum by \(q(r)\), giving
\[
  \sum_xq(x)=\sum_x q(x+r)=q(r)\sum_xq(x).
\]
 Passing to the quotient by the
radical therefore changes the sum by a positive integer.

For completeness, let \(q\) be a nondegenerate homogeneous
quadratic character on a finite abelian group \(B\). Writing \(b(x,y)=q(x+y)/(q(x)q(y))\), character orthogonality gives
\[
  \left|\sum_xq(x)\right|^2=\sum_h q(h)\sum_y b(h,y)=|B|.
\]
Denote this Gauss sum by \(G\). To see that its
phase is an eighth root of unity, put \(N=|B|\) and choose
four integer squares summing to \(N^2-1\). Left quaternion
multiplication by this vector gives an integer similarity
\(Q\) of the sum of four squares with
\(Q^{t}Q=(N^2-1)I\), invertible on \(B^4\).
Expanding the quadratic character and its symmetric
bicharacter gives

\[
  \prod_i q((Qv)_i)=(\prod_i q(v_i))^{N^2-1}
=\overline{\prod_i q(v_i)}.
\]
 Here \(q(Nx)=q(x)^{N^2}=1\).
Thus \(G^4\) is real. This also covers \(N=1\).
As \(|\xi|=1\), the triple comparison follows.

An additional determinant occurs for a change of variables.
Let \(L=S\) and take the rational unitary \(\delta\) of
Lemma~\ref{lem:cm-standard-local-isometry}. Put
\(A=1-\delta^{-1}\), \(I=L\cap\delta L\) and \(j=[L:I]\).
Then for \(F\) on the theta line of \(L\),
\begin{equation}\label{eq:cayley-multiplier}
   F\circ\delta^{-1}=c_\delta T_{L,\delta L}F,\qquad
              c_\delta^8=(\det\delta)^4 .
\end{equation}

To compute the scalar, periodize the Gaussian kernel in its
first variable:
\[
 P(v,z)=\sum_{x\in L}\alpha(x)
      U_x\big(e^{\pi H(\cdot,z)}\big)(v).
\]
Positive definiteness gives normal convergence on compact sets
by a Euclidean lattice Gaussian, and \(P(\cdot,z)\) lies in the
theta line. Integrate \eqref{eq:cayley-multiplier}'s first
identity for this section, evaluated at \(z\), against
\(e^{-\pi H(z,z)}\,dz\) on a fundamental cell of \(I\).
Before inserting \(c_\delta\), the right side integrates to
\(\operatorname{vol}(I)/\sqrt j\), where \(\operatorname{vol}\)
of a lattice denotes covolume. Indeed in the expansion of
\(T P\) the oscillatory factors are characters
\(e^{2\pi i E(x+y,z)}\) for \(x\in L\), \(y\in\delta L/I\);
integration leaves just \(x+y=0\).

On the left the summand for \(x\) is
\[
 q_A(x)\,e^{-\pi H(A(z-A^{-1}x),\,z-A^{-1}x)},\qquad
 q_A(x)=\alpha(x)\exp[-\pi H((A^{-1}-\tfrac12)x,x)] .
\]
This is completion of the square using
\(A+A^*=A^*A\). Its absolute Gaussian has exponent
\(-\pi H(Az,Az)/2\) after recentering. The polarization
of \(q_A\) on \(L\) is \(e^{-2\pi i E(A^{-1}x,y)}\).
Its radical consists precisely of \(AI\): self-duality
says \(A^{-1}x\in L\), and \(x\in L\) forces
\(\delta^{-1}A^{-1}x\in L\). For \(x=Av\) with \(v\in I\),
substituting \(x=v-\delta^{-1}v\) and

\[
  \alpha(\delta^{-1}v)=\alpha(v)
\]
 into the semicharacter
law gives \(q_A(x)=1\). Hence it descends to a
nondegenerate homogeneous quadratic character on \(L/AI\).
Group the sum by this quotient and unfold the translates by
\(I\). Absolute integrability follows from the finite number
of recentered Gaussians. Diagonalizing \(A\) in unitary
coordinates yields
\[
 \int e^{-\pi H(Az,z)}\,dz=
      (\det A)^{-1}\int e^{-\pi H(z,z)}\,dz
       =\operatorname{vol}(L)/\det A .
\]
Each eigenvalue has positive real part since \(1\) is excluded
from the spectrum of \(\delta\). The last equality follows in
orthonormal coordinates from the symplectic covolume one of a
self-dual lattice. Since \([L:AI]=j|\det A|^2\), division by
the right-hand integral gives
\[
  c_\delta=\frac{|\det A|}{\det A}
      \frac{1}{\sqrt{[L:AI]}}\sum_{x\in L/AI}q_A(x).
\]
The normalized sum has eighth power one.
Finally \(A^*=-\delta A\), so
\(\overline{\det A}=(-1)^m\det\delta\det A\);
this proves the asserted eighth power with its sign and exponent.
\end{proof}

\paragraph{Returning to the local coefficients.}
Apply Lemma~\ref{lem:theta-cayley-scalar} before taking the
fourth tensor and returning to the local argument lattice.
Set \(K_i=(T_{S,L_i}F_0)\circ g_i\). Transport of \(U_x\)
under \(\delta\) and the triple comparison give
\(K_0=cK_1\) with \(c^8=(\det\delta)^4\). Explicitly, in
the \(g_1\) coordinates the transported section is
\(T_{\delta S,L_1}(F_0\circ\delta^{-1})\).
Use \eqref{eq:cayley-multiplier} and then compose the two
intertwiners. Since the theta line of \(M^4\) is the fourth
tensor power of the line of \(M\), there exists a theta
vector \(\Theta\), unique up to a fourth root of unity,
with \(\Theta^{\otimes4}=\kappa K_0\). This fixes its
complex norm at every conjugate of the determinant scalar.

For each finite embedding 
\[
  \Theta^{\otimes4}=\kappa K_0=(\kappa c)K_1
\]
and \(|\kappa c|_p=1\) by \eqref{eq:kappa-determinant}.
Every summand of \(T_{S,L_1}F_0\) is a product of
one-variable translates with prime-to-\(p\) denominators,
and its normalization is a sign times a unit square root
of an inverse index. A further translation by a point of
order prime to \(p\) adds points of the same kind in the
standard coordinates: a prime-to-\(p\) multiple of its
\(g_1\)-image lies in \(L_1\), and a prime-to-\(p\) index
carries that lattice into \(S\).

Since the two translation arguments belong to the rational space of \(S\), write \(E(y,y')=a/b\) with \(b>0\). The composition phase satisfies
\[
  \bigl(e^{\pi iE(y,y')}\bigr)^{2b}=e^{2\pi ia}=1.
\]
This gives the normalized orthogonal fourth tensor expansion
used in Section~\ref{subsec:cm-models}. Section~\ref{subsec:baseline-jets} expresses the
coordinate change on the argument lattices by CM matrices
integral at the place and deduces integrality of derivatives.
The positive gains also use the incident chart and its
linear forms in Section~\ref{subsec:incident-sections}.
